\chapter{Religion, Doubt, and Secularization}
\section{A Charm for Passing Every Examination}
First, pick up something small: the protective charm hanging from a schoolbag's zipper.

Perhaps you know a classmate like this. Before examination week, a little embroidered pouch promising success appears on her bag, or an amulet brought back from Sensoji, the Confucius Temple, or Yonghe Temple. Ask whether she believes in it and she will probably laugh: "Just for good luck." Then she turns over her phone case, revealing a paper talisman folded into a tiny square inside.

That answer is interesting. Ask whether she believes in gravity and she will not say, "Just for good luck." Ask whether humans can photosynthesize and she will say no. Yet with this charm, her answer is neither belief nor disbelief but "It can't hurt to wear it."

Notice her position: she does what a believer might do---seeking a charm, wearing it, touching it before examinations---but refuses the identity of believer. Behavior and belief occupy separate ledgers. Nor is this merely her personal oddity. It is widespread across humanity: Japanese people register births at shrines, marry in white dresses in churches, and invite monks to chant at funerals; Europeans attend church less but still celebrate Christmas; your elders may claim to believe in nothing while visiting ancestral graves at Qingming, putting up the character for good fortune at New Year, and sharing lucky-carp pictures faster than anyone.

Is the apparently simple question "Is this person religious?" really so simple? Is that charm belief, superstition, tradition, a placebo, or a social signal? If someone can act without believing, what does belief mean?

We begin with this charm to explore one of intellectual history's greatest tensions: do gods exist? Can humans know the answer? What would happen if most people ceased believing? And a question you may not have considered: is disbelief itself a right requiring protection?

\section{Falling Blossoms at the Western Residence}
Now enter a scene with me. Time: around the late fifth century CE, more than fifteen hundred years ago. Place: the Western Residence of Xiao Ziliang, Prince of Jingling, in Jiankang, capital of the Southern Qi dynasty, today's Nanjing.

The Western Residence was the country's liveliest cultural salon. Xiao Ziliang, an imperial prince and devout Buddhist, supported eminent monks and attracted scholars from across the land. Imagine costly incense rising straight from a burner; monks' robes crowded among scholars' broad garments; discussions of poetry and Buddhist doctrine, not loud, but with every word weighed by listeners. Admission to this room was itself a scholar's career prospect.

Among the guests was Fan Zhen, a scholar of modest rank but considerable reputation for speaking boldly. That day Xiao Ziliang challenged him. Historical records preserve their exchange. Xiao Ziliang asked:

\begin{quote}
If you do not believe in karmic consequences, how can the world contain wealth and honor, poverty and low status?
\end{quote}
In other words: if Fan Zhen rejects karmic recompense, why are some born rich and powerful and others poor and humble? If these differences are not the fruits of good and evil in previous lives, where do they come from?

This was not casual conversation but a check in a chess game. Within the prevailing framework, karmic recompense was the master key to social differences: prosperity rewarded merit from a previous life; suffering repaid previous karma. The explanation did two jobs at once, explaining the world and stabilizing order. The poor accepted their lot, the rich felt reassured, goodness earned interest, and wrongdoing incurred a bill. If Fan Zhen failed to answer well, what would be dismantled was not merely one theological opinion but an entire scaffold for explaining the world.

The Book of Liang, in its biography of Fan Zhen, records his answer:

\begin{quote}
Human birth resembles blossoms on one tree, emerging from one branch and opening together upon one stem. Falling with the wind, some brush past curtains and land upon fine mats; others pass beyond fences and fall beside the privy.
\end{quote}
Put plainly, people resemble flowers growing and blooming together until the wind scatters them. Some petals drift through curtains onto luxurious cushions, others over fences into a cesspit. The flower on the cushion is Your Highness; the one in the cesspit is me. Wealth and poverty indeed follow different paths---but where is karmic recompense in that?

What makes this powerful? Fan Zhen neither insults Buddhism nor attacks Xiao Ziliang, nor even raises his voice. He simply changes keys: chance explains the difference without a previous life's account book. Flowers bloom and fall; where the wind takes them is fortune, not a bill. This was explosive. If status is merely blossoms on the wind, the assertion that people deserve all their suffering collapses. Yet so does the hope that goodness will be rewarded, a hope helping countless people endure and remain good.

What followed reveals the stakes. Historical accounts say Xiao Ziliang sent the eloquent Wang Rong to persuade Fan Zhen. His offer amounted to this: your claim that the spirit perishes is unsound, yet you cling to it; do you not fear harming established moral teaching? With your talent, becoming a palace secretary would be easy. Why deliberately oppose everyone? Destroy the essay. Fan Zhen laughed. If he were willing to trade his arguments for office, he replied in effect, he would long since have obtained a much higher post than that.

Remember "trading arguments for office." A scholar, pressed simultaneously by imperial power, religion, and career prospects, refused to withdraw a view he considered true. This is doubt in its historical setting: not always calm calculation in a laboratory, but sometimes laughter in a drawing room, followed by the bill for laughing.

\section{Is Doubt Dangerous, or Clear-Sighted?}
Bring the Western Residence's conflict into full light before declaring a winner.

What exactly worried Xiao Ziliang's side? On the surface, the dispute was metaphysical: does the spirit survive death? Beneath it ran three thicker cords.

The first was explanation. People need an account of the world's unevenness: why is one person born in a palace and another beside a cesspit? Karmic recompense may make you uncomfortable today, but it offers an answer. Fan Zhen replaces it with chance. Intellectually this may be more honest, yet emotionally it resembles a doctor announcing that your illness has no cause. Some feel liberated; others feel chilled.

The second was morality. Without rewards for goodness, punishment for evil, or an afterlife's court of judgment, why be good? Skeptics in every age face this question. Xiao Ziliang's guests asked Fan Zhen; religious elders ask nonreligious young people at dinner today. It deserves serious attention, because it is not mere provocation. Many societies really have hung moral education on religion's nail. Whether pulling the nail out will bring down the plaster is a genuine concern.

The third was power. Xiao Ziliang was not only a believer but a ruler. Karmic recompense was exceptionally useful to rulers, translating present inequality into individuals' moral grades from previous lives. Complaining about society became questioning cosmic justice. Fan Zhen's blossom analogy unsettled the powerful because it burned the ledger. If wealth and poverty are merely outcomes of the wind, a dangerous question follows: why must I accept that the wind placed you on the cushion? Doubt here is not merely philosophy; it can unsettle order.

But Fan Zhen's side carries equal weight. The skeptic holds honesty: should an unsupported explanation be accepted simply because it is useful, comforting, or convenient for government? "This comforts people" and "This is true" are different statements. Someone may believe out of kindness, but kindness cannot substitute for evidence. Nor does doubt automatically demolish everything. Fan Zhen himself is remembered for filial devotion and integrity; he rejected an afterlife and still lived morally. Here lies our real question, larger than whether gods exist:

\textbf{After an ancient explanation is questioned, can the work of explaining the world, sustaining morality, and organizing society still be done? Who will do it?}

This is not merely a fifteen-hundred-year-old problem. It plays out today in every family where someone announces, "I no longer believe."

\section[Disbelief and the Case for Belief]{Four Kinds of Disbelief, and a Full Hearing for Belief}
At least half the arguments about religion arise because people use different terms as if they meant the same thing. Let us distinguish four common kinds of "disbelief," each with an everyday example.

\textbf{Atheism}: asserting that gods do not exist. This is not indifference or lack of interest but a definite judgment, like saying there are no round squares. Fan Zhen was one atheist in ancient China. His argument was not that the afterlife did not interest him but that mind could not exist apart from body, leaving no soul to receive recompense after death. An older example prevents us from imagining atheism a modern Western invention. Ancient India's Charvaka school, roughly contemporary with Shakyamuni or earlier, held that the world consisted of earth, water, fire, and air; bodily death ended the person, with neither afterlife nor karmic recompense. Most of its writings are lost. We know it largely through Buddhist and Jain attacks: opponents' filing cabinets preserved its name.

\textbf{Agnosticism}: asserting neither existence nor nonexistence, but saying that one does not know and perhaps cannot know. The nineteenth-century British naturalist Thomas Huxley coined the term. His position resembles a strict judge: the prosecution asserts God exists and the defense asserts the opposite; neither provides enough evidence for a verdict, so judgment must remain open. This is fidelity to standards of evidence, not compromise for its own sake. Imagine never having visited Iceland and being asked the color of the sign on the third shop along a particular street. The honest answer is not a guess but "I don't know." Agnostics regard much evidence about gods as being in that condition.

\textbf{Deism}: belief in a creator who no longer intervenes after making the world. Popular during the eighteenth-century European Enlightenment, it uses the image of a clockmaker: God constructs the precise cosmic clock, winds it, and leaves its gears to turn without answering prayers or sending miracles. Thinkers such as Voltaire broadly held this position. A school analogy would be a teacher who sets all classroom rules on the first day and never returns, leaving the class to run by them. Deism removes whole questions about effective prayer and authentic miracles: God exists, but is permanently on holiday.

\textbf{Secular humanism}: this changes the subject altogether. Rather than first answering whether God exists, it asks how humans should live. Morality needs no divine approval; compassion, reason, and concern for humanity's shared fate can sustain a benevolent, meaningful life. Its core is not "There is no God, so anything goes," but "Without divine commands, people must take greater responsibility for their choices."

Together these positions reveal an easily missed distinction. Atheism answers an existence question; agnosticism a knowledge question; deism an intervention question; secular humanism a question about living. They are not four points on one scale but answers to four different questions.

\textbf{Now a necessary exercise: state the believer's case in full.} Narratives favoring skepticism often reduce faith to ignorance, cowardice, or needing a crutch. That is lazy. Listen seriously: for many believers, faith is not signing a proposition saying God exists but a relationship and a way of life. It provides someone to pray to amid disaster, a story before death extending beyond decay, and a place among fellow worshippers where one need not explain who one is. Religion also bundles meaning---why I live; community---who my people are; ritual---how birth, adulthood, marriage, and death are solemnly marked; and moral grounding---goodness as a requirement with a source, not mere taste. You may accept none of this. But unless you can restate it in terms believers recognize, you have not understood what you intend to refute. This book repeatedly uses that principle: complete the other person's argument before replying. It is the admission ticket to serious disagreement.

\section[Thinking Bridge: Four Questions]{Thinking Bridge: Split Belief into Four Questions}
Here is a portable tool. When religion next produces an argument at dinner, in a group chat, or in the news, do not immediately take sides. Separate the questions first.

\textbf{Question one: existence. Do gods, spirits, karmic recompense, or an afterlife exist?} This is a factual question, with in principle one truth regardless of how many opinions conflict. Fan Zhen and Xiao Ziliang were debating this question.

\textbf{Question two: knowledge. Can humans know about gods?} Notice its independence from the first. One can believe in God while admitting inability to prove it; many believers describe faith precisely as trust where evidence is insufficient. Or one can regard the question as beyond human cognitive limits, as agnostics do. What an answer is and whether I can know it occupy different floors.

\textbf{Question three: ethics. Does being good require God?} This too is independent. History contains morally rigorous atheists, including Fan Zhen, and people committing evil in God's name. Both sets of examples show that belief and good conduct are not linked by one simple line. This does not mean religion has no moral effects. It means their existence, size, and value require particular examination, not a verdict that without religion people become lawless or that all religion is a moral shackle.

\textbf{Question four: politics. What role should religion have in public life?} Where should boundaries lie between churches and schools, states and temples, law and doctrine? This is still further from personal belief. Devout believers have advocated separating religion and state because binding them together corrupts both. Atheists have supported religious groups' public roles because of their usefulness in disaster relief and education.

The tool's use fits one sentence: when an argument starts, ask which question is being debated. Much talking past one another occurs when one person answers the first and another the fourth. "God does not exist" becomes "Churches should close" in the listener's ears; "Schools should not proselytize" becomes "Your faith is false." Once the questions are separated, many arguments stop because the parties realize they were discussing different matters. Genuine disagreements at least reach the correct table.

\section{Two Doubts Written on Paper}
Now read a short primary source. Questioning religious explanations is no modern monopoly; it has a written record.

First, Fan Zhen. After the Western Residence debate he set out his position in On the Extinction of the Spirit. When Emperor Wu of Liang took the throne, he personally issued a rebuttal and organized dozens of nobles, officials, and eminent monks against it. Some responses from more than sixty opponents survive. The essay's most famous argument reads:

\begin{quote}
Spirit relates to substance as sharpness to a blade; bodily form relates to function as a blade to sharpness. The name sharpness does not mean blade, nor blade sharpness. Yet apart from sharpness there is no blade, and apart from the blade no sharpness. Never has sharpness survived the blade's destruction. How could spirit remain when bodily form has perished?
\end{quote}
(On the Extinction of the Spirit, preserved in the Hongming ji, a Buddhist collection compiled by Sengyou during the Southern Dynasties)

In ordinary language, mind and body relate as sharpness and blade do. They are distinct names referring to different things. But without the blade, where could sharpness be? We never hear of a vanished knife whose sharpness remains. Why suppose the mind remains after the body's death?

The analogy's force is its directness. Sharpness is not a tiny object hidden in a knife but its property or function. Likewise, mind is not a little lodger inside the body but bodily activity itself. Lodgers can move; properties cannot. A flame cannot exist separately from burning, or sharpness from a blade. Why should mind exist separately from body?

The essay's survival is also worth noting. It lived in its opponents' filing cabinet. Fan Zhen's collected writings did not survive intact, but Buddhist compilers included the entire essay in the Hongming ji so successive generations of monks could learn to refute it. Many heterodox voices survive this way: enemies copied their words precisely in order to criticize them. This is itself a lesson in sources. Which ancient books we can read often depends on who once considered them important, whether worthy of veneration or concerted attack.

Now an earlier source, with a very different tone. In the Analects, "Those Who Advanced," Zilu asks Confucius how to serve spirits. Confucius replies:

\begin{quote}
Before you can serve people, how can you serve spirits?
\end{quote}
Zilu ventures another question, about death. Confucius says:

\begin{quote}
Before you understand life, how can you understand death?
\end{quote}
Together: before attending properly to the living, why discuss spirits? Before understanding life, why discuss death?

Notice the stance. Confucius neither denies existence like Fan Zhen nor affirms it like a priest. He \textbf{sets the question aside}: it need not be answered now because other matters take priority. The Analects, "Yong Ye," adds a related statement: attending to people's proper duties and respecting spirits while keeping them at a distance may be called wisdom. Respect, but do not draw near.

Side by side, these sources show two ancient forms of doubt. Fan Zhen's is \textbf{denial}, arguing that the claim fails. Confucius's is \textbf{suspension}, judging neither yes nor no and returning attention to human affairs. Both positions were already occupied in China fifteen hundred years ago. Belief and disbelief are therefore not modern imports but positions humans arrive at once they begin thinking, and each allows different ways of sitting.

\section[Science and Religion: Always at War?]{Have Science and Religion Really Fought for Centuries?}
Next compare two explanations of the same facts: what relationship has science had with religion since its rise?

\textbf{Explanation one: the warfare model.} This is the most widespread version, readily found online. Religion represents ignorance and authority; science represents reason and freedom; they are destined enemies, and every scientific advance must cross the church's line of fire. Familiar characters include Bruno burned at the stake, Galileo tried, and churches obstructing evolution. In the nineteenth-century English-speaking world, influential books wrote the whole history of science as conflict with religion. Warfare became the public's default image.

There is truth here. Conflicts genuinely occurred. In 1633, the Inquisition tried Galileo and forced him publicly to renounce Earth's movement around the Sun; he spent his remaining years under house arrest. Copernicus's book did appear on the church's Index. The explanation identifies real oppression. Its limitation is excessive neatness: it cuts centuries into a good-versus-evil film, removing scenes that do not fit.

\textbf{Explanation two: the complex-relationship model.} Most historians of science today hold this view: relations vary with period, question, and people, with abundant examples of conflict, cooperation, and mutual irrelevance. Consider the scenes cut by the first explanation:

Copernicus himself was a church cleric. His heliocentric book appeared in 1543 dedicated to the pope and was not prohibited on publication; only more than seventy years later, in 1616, was it listed pending correction. Mendel, discoverer of the laws of heredity, was an Augustinian friar who experimented with peas in a monastery garden. Lemaitre, whose proposal that the universe began in a primeval atom preceded Big Bang theory, was a Belgian Catholic priest. The Vatican not only refrained from suppressing it but briefly wanted to enlist it in support of creation. Lemaitre himself advised against this: scientific theories and faith should not be bound together. Outside Europe, practical requirements to determine prayer direction and Ramadan's timing directly stimulated astronomy and mathematics during the medieval Islamic world's golden age. Algebra itself takes its name from Arabic, and many astronomers also held mosque positions. Europe's earliest universities grew almost entirely from church schools.

Even Galileo's trial, the warfare model's leading exhibit, is not a pure specimen. Galileo remained Catholic throughout life. The trial entangled authority over biblical interpretation, factional struggles inside the church, and powerful people's wounded pride at his written mockery. Scientific conclusions were only one fuse. This is an easily misread counterexample: the event most readily treated as a standard image of science fighting religion resists being read as pure science against pure religion.

Compare the explanations. Warfare captures real oppression at the cost of oversimplifying history and portraying both religion and science as internally unified camps. In reality, both contain arguments. Monks can be scientists, believers can support new theories, and scientists' quarrels with one another can be as fierce as their quarrels with churches. Complexity fits the evidence better but sacrifices a satisfying story and a quick answer about who the good people are. You need not decide today that one explanation is entirely correct. Take away a sense of proportion: \textbf{it is wrong to say science and religion have always fought, and wrong to say they never conflict. The answer depends on the period, the issue, and the people involved}.

\section[Further Challenge: Secularization]{Further Challenge: What Does Secularization Mean?}
Now introduce our theoretical equipment. Secularization appears constantly in public discussion, with thoroughly tangled meanings. Sociologists have studied it for more than a century, beginning with the discovery that it refers to at least three different things. Mixing them produces many arguments.

\textbf{First layer: institutional separation.} State and religious institutions part company: government establishes no official church, churches cannot directly command government, and public schools do not proselytize. This is the hardest layer, written into constitutions and law. France's 1905 law established a radical separation, removing even religious signs from public institutions. The U.S. Constitution approaches the matter differently: Congress may neither establish religion nor prohibit its free exercise. Countries dispute the merits and extent of this secularization differently, but notice that it concerns \textbf{institutional arrangements}, not anyone's inner faith. A state separating religion and government can be full of devout people; the United States is a ready example.

\textbf{Second layer: declining belief.} Fewer people attend church or identify as religious, and prayer and worship become less frequent. This concerns \textbf{minds and behavior}. Twentieth-century Western Europe is the classic case: churches remain, but Sunday seats empty. Sociologists once regarded this as a one-way street: modernization would shrink faith until religion withered like an appendix.

\textbf{Third layer: changing forms.} Religion does not vanish but changes how it exists, from everyone's default operating system to an optional product on a shelf, from something bound to community, family, and identity to a spiritual resource individuals choose as needed. Scholars studying Europe summarized this as "believing without belonging": retaining some transcendent conviction without joining churches or following their rules. In Japan, most people call themselves nonreligious yet seek shrine blessings for newborns, church weddings, and Buddhist funerals. Religion functions less as belief than as customary procedures for life's stages. Return to the examination charm: it perfectly exemplifies this third layer, with ritual retained, belief quietly removed, and comfort, tradition, and a little "just in case" left behind.

Only after separating these layers can we reach the chapter's deepest theory. German sociologist Max Weber used the famous phrase \textbf{the disenchantment of the world}. He meant more than declining belief in gods: the world's whole character changes. Before disenchantment, intentions lie behind everything. Thunder expresses Heaven's attitude, illness follows an offense, and harvests reflect favor. The world resembles a forest full of spirits and wills, dangerous but meaningful, with beings to address everywhere. After disenchantment, it becomes a calculable causal machine: thunder is atmospheric discharge, illness viruses and cells, harvests seeds, fertilizer, and climate. It becomes transparent and controllable, but also silent. Seeking a hospital before making ritual amends when ill is disenchantment embodied in your own life. Weber's insight was that rationalization brings enormous power while taking something away: meaning no longer grows spontaneously from the world, and people must make it themselves. Disenchantment is no triumphal march of truth defeating superstition, but a transaction involving gain and loss.

The theory is ready; now comes the challenge. Mid-twentieth-century sociologists confidently predicted religion's disappearance with modernization. What happened? The prophecy came true only halfway. Western European churches did empty, yet in many other places religion persisted, grew more fervent, or changed form. New spiritual consumption---meditation, astrology, lucky carp, mind-body-spirit courses---filled spaces vacated by traditional churches. Peter Berger, an American sociologist and once a leading secularization theorist, later publicly acknowledged that the prediction had been wrong: the world remained as intensely religious as before, sometimes more so. Canadian philosopher Charles Taylor offered a subtler account in A Secular Age: modernity's decisive shift is not from belief to disbelief but \textbf{from belief as the default to belief as a choice}. Centuries ago, disbelief was almost unimaginable, with God as unquestioned a background as gravity. Today belief and disbelief are options on a shelf. Everyone must select and carry their own faith, which is therefore less secure than their ancestors'.

The further challenge is thus not a conclusion but a harder question. If secularization means religion changing from air into an option, are people today freer than their ancestors, or lonelier?

\section[Today: Faith and Public Rights]{At Today's School Gate: Headscarves, Lucky Carp, and Constitutional Rights}
These old questions continue in today's news and ordinary life. Consider three scenes.

\textbf{Scene one: headscarves at French school gates.} In 2004, France prohibited conspicuous religious symbols among pupils in public primary and secondary schools, including Muslim girls' headscarves, oversized crosses, and Jewish skullcaps. Our tool identifies a question-four dispute: not about the scarf itself, but where separation of religion and state should lie. Supporters argue that public schools are neutral state spaces, where children should leave religious identity outside to avoid peer and community coercion. Some girls may not wear scarves voluntarily; school neutrality protects their space to refuse. Opponents argue that wearing a scarf is personal religious expression; the state removing it in neutrality's name violates freedom. Moreover, the practical burden falls almost entirely on Muslim pupils, placing special demands on one group behind neutrality's mask. Both sides invoke freedom and protection while reaching opposite conclusions. That is public disagreement's real difficulty, and why separating questions matters more than taking sides.

\textbf{Scene two: secularism in India's Constitution.} Secularization has no single worldwide meaning. India's Constitution declares a secular state, yet its version is almost the opposite of France's: rather than excluding religion from public life, the state maintains equal distance from religions. It attends to Hindu festivals and also Muslim, Sikh, and Christian ones; a Hindu majority does not make it a Hindu state. The same word means separation in France and care for all in India. Before debating religion's public role, ask which secularism the speaker means.

\textbf{Scene three: should religion speak on public issues?} Every society debates this today. Let us fully state both cases, our second steelmanning exercise. The withdrawal case says public debate requires reasons everyone can examine. "My scripture demands this" offers nonreaders a conclusion, not a reason. Historically, religious public power has often harmed minority believers and nonbelievers first. Public life should therefore address people as citizens, not adherents. The participation case responds that some of history's most just movements used religious language. Martin Luther King Jr. was a Baptist minister, Black churches organized the U.S. civil rights movement, and biblical images filled his speeches. Gandhi's independence campaigning drew on Hindu traditions of asceticism and compassion. Religious groups still operate many disaster shelters, soup kitchens, and remote schools. Asking believers to leave faith at home when entering public life splits them in two, because faith often motivates their public concern. Both hold real truths. A more honest compromise allows religious reasons into public discussion while subjecting them to public rules: faith may motivate opposition, but persuading those who do not share it requires reasons they too can understand.

Finally, a right directly relevant to you. Article 18 of the United Nations' 1948 Universal Declaration of Human Rights states that everyone has freedom of thought, conscience, and religion, explicitly including freedom to \textbf{change} religion or belief. Pay particular attention to change. Historically, religious freedom's first stopping point was permission to follow one of several approved faiths. Its fuller form contains three freedoms: to choose a faith, to move from one to another, and to believe in none. Article 36 of China's current Constitution grants citizens freedom of religious belief, encompassing both belief and nonbelief and prohibiting coercion either way. These lines were not always there. For much of human history, conversion and disbelief were grave offenses. Fan Zhen faced an emperor's organized attacks merely for an essay, refusing through laughter to trade arguments for office. Even today, in some places openly leaving an inherited religion risks family rupture or personal safety. Freedom not to believe is no self-evident atmosphere. People fought, wrote, and argued to put it into texts, and it still needs protection. Here is the other side of the examination charm: the classmate can half-believe, hang it up or remove it, and jokingly say "just in case" because she happens to occupy a position permitting belief, doubt, or neither. That position did not come free.

\section{Your Question: Two Readings of a Paper Charm}
Return to the schoolbag.

You now have two complete readings of its examination charm. The first sees harmless comfort, tradition's lingering warmth, and an outlet for anxiety. Its wearer deceives nobody, consistently admitting it is for good luck. Humans have always needed ritual, whose value need not pass a belief checkpoint. The second sees a small dishonesty: facing something important, she borrows what she intellectually rejects, speculating on "just in case." If everyone summons gods when useful and shelves them afterward, what weight remains in the word conviction?

Enlarge the question and we reach the chapter's foundation: \textbf{if most people in a society cease believing in gods, will morality collapse, or acquire a new foundation?}

Those expecting collapse point to karmic recompense's abrupt departure, Weber's loss of meaning, wrongdoing's reduced costs without an afterlife, and the loneliness left when religious communities dissolve. Those expecting new foundations point to Fan Zhen's integrity without belief in an afterlife, secular humanists' commitment to kindness without gods, and the high trust sustained in Nordic societies with empty churches. They also ask: if goodness itself requires a heavenly ledger, is that morality not already fragile?

Both sides have evidence and unanswered questions. What is your answer, and why? Whichever you choose, you enact the chapter's final fact: when belief becomes a choice, believers and nonbelievers alike must carry their own reasons. Nobody else can sign for them.
